\documentclass[
sn-mathphys,
Numbered,
onecolumn]
{sn-jnl}% Default with double column layout

\usepackage{graphicx}% Include figure files
\usepackage[title]{appendix}%
\usepackage{dcolumn}% Align table columns on decimal point
\usepackage{manyfoot}%
\usepackage{graphicx}
\usepackage{currfile}
%\usepackage{dcolumn}% 
\usepackage{multirow}
\usepackage{chemformula}
\usepackage{bm,soul}
\usepackage{color}
%\usepackage{hyperref}
%\usepackage[all]{hypcap}
\usepackage{gensymb}
\usepackage{amsmath}
\usepackage{amssymb}
\usepackage{dcolumn}
\usepackage{float}
\usepackage[export]{adjustbox}
\usepackage{multirow}
\usepackage{rotating}

\newcommand{\Ang}{\AA$^{-1}$}
\newcommand{\etal}{\textit{et al.}}
\newcommand{\NI}{\noindent}
\newcommand{\NL}{$\newline$}
\newcommand{\ueV}{$\mu$eV}
\newcommand{\XE}{\ensuremath{\chi''(E)}}
\newcommand{\XQE}{\ensuremath{\chi''(Q,E)}}
\newcommand{\SQE}{\ensuremath{S(Q,E)}}
\newcommand{\minus}{\scalebox{0.75}[1.0]{$-$}}
\newcommand{\soe}{$S_{Q}(E)$}
\newcommand{\soq}{$S_{E}(Q)$}
\newcommand{\xoq}{$\chi''_{E}(Q)$}
\newcommand{\xoe}{$\chi''_{Q}(E)$}

\definecolor{grey}{gray}{.5}%normally{.6}
\newcommand{\grey}[1]{{\color{grey}{#1}}}
\newcommand{\red}[1]{{\color{red} #1}}
\newcommand{\blue}[1]{{\color{blue} #1}}
\newcommand{\magenta}[1]{{\color{magenta} #1}}
\definecolor{malacol}{rgb}{0.12,0.69,0.44}
\newcommand{\mala}[1]{{\color{malacol}{#1}}}

\newcommand{\uSR}{$\mu$SR}
\newcommand{\uB}{\ensuremath{\mu_\textrm{B}}}
\newcommand{\herb}{ZnCu$_3$(OH)$_{6}$Cl$_{2}$}
\newcommand{\um}{$\mu$m}

\newcommand{\ave}{Cu$_5$(VO$_4$)$_{2}$O$_2$CsCl}
\newcommand{\znave}{ZnCu$_4$(VO$_4$)$_{2}$O$_2$CsCl}
\newcommand{\zntwoave}{Zn$_2$Cu$_3$(VO$_4$)$_{2}$O$_2$CsCl}
\newcommand{\aveser}{Zn$_x$Cu$_{5-x}$(VO$_4$)$_{2}$O$_2$CsCl}

\renewcommand*{\thesection}{\Roman{section}}
\renewcommand{\thefigure}{S\arabic{figure}}
\renewcommand{\thetable}{S\arabic{table}}

% % %   SUPPLEMENTAL   % % % 
\begin{document}

\begin{center}
\textbf{\large Supplemental Material for ``Magnetic ground states and excitations in Zn-doped averieite - a family of oxide-based $S=1/2$ kagome antiferromagnets"}
\end{center}
%\appendix

\section{Crystal structures}
\subsection{Averievite, \ave}

Neutron diffraction patterns of averievite collected on HRPD at temperatures between 1.5 and 100~K are shown in Fig.~\ref{FigCu5HRPD} and the Rietveld refinement at $T=1.5$~K (Fig.~\ref{FigCu5HRPDRietveld}) with the resulting lattice parameters and atomic positions displayed in Table~\ref{avv_1K_params}. Important Cu-Cu distances and superexchange angles are given in Table~\ref{TableAveAngles}.


% Figure Cu5 HRPD temperature scans
\begin{figure*}[!ht]
\centering
\includegraphics[width=0.99\textwidth]{Cu5_hrpd_all_T.png}
\caption{Averievite, HRPD neutron diffraction data at 1.5, 25, 50, 75 and 100~K. The data are presented as the square root of the intensity to facilitate comparison of the data sets at higher $Q$ where the peak intensity is lower. No new peaks or peak splittings are seen as temperature decreases, indicating no further crystallographic transition below 100~K.
}
\label{FigCu5HRPD}
\end{figure*}
% Source: Madeleine's thesis

%The $T=100$~K 11-BM diffraction pattern was fitted with a Rietveld refinement of averievite in space group $P12_1/c1$ with lattice parameters $a$, $b$, $c$ and $\beta$, and a Pawley fit with lattice parameters $a^\prime = a$, $b^\prime = 3b$, $c^\prime = c$ and $\beta^\prime = \beta$. 
%A CuO impurity phase was also included. The profile parameters for the two averievite phases were set to be equal and the refinement is shown in Fig.~\ref{FigCu511BM100K}.

% Figure Rietveld refinement Cu5 HRPD 1.5 K
\begin{figure*}[!h]
\centering
\includegraphics[width=0.99\textwidth]{Cu5_hrpd_1K.png}
\caption{Averievite, \ave, HRPD neutron diffraction data collected at 1.5~K (black) with a Rietveld refinement in red and the difference plot in grey. The refinement showed a $\sim1$\% CuO impurity in the sample. The peak positions of the main averievite phase and of CuO are shown by green and blue ticks, respectively. The excluded regions are Bragg peaks from the aluminium sample can.
}
\label{FigCu5HRPDRietveld}
\end{figure*}
% Source: Madeleine's thesis

\begin{table}[ht!]
    \centering
    \begin{adjustbox}{max width=\textwidth}
    \begin{tabular}{c c c c c c c } \hline \hline
    %\multicolumn{7}{c}{Atomic parameters} \\ \hline
     Atom  & Wyckoff site & $x$ & $y$ & $z$ & Occupation & $U_\mathrm{iso}$ (\AA$^2$) \\ \hline
  Cu$_{k1}$ & $2c$ & 0 & 0 & 0.5  & 1 & 0.00115(43)  \\
  Cu$_{k2}$ & $4e$ & 0 & 0.8005(4) & 0.2288(3)  & 1 & 0.00115(43)  \\
  Cu$_h$ &  $4e$ & 0.7296(4)  & 0.5334(5) &  0.1668(3)  & 1 & 0.00415(55) \\
 Cl &  $2a$ &  0 &  0 &  0 & 1 & 0.0177(8) \\
 Cs &  $4e$ &  0.4242(7) & 0 & 0 & 0.5  & 0.0120(16) \\
 O$_1$ & $4e$ & 0.7610(4)  & 0.6732(5) &  0.7418(3)  & 1 & 0.00174(31)  \\
 O$_2$ & $4e$ & 0.7640(4)  & 0.2663(5) & 0.2700(4)  & 1 & 0.00174(31)  \\
 O$_{3}$ & $4e$ & 0.7626(4)  & 0.4958(5) & \minus 0.02130(31)  & 1 & 0.00174(31)  \\
 O$_{Cs}$ & $4e$ & 0.4950(4) & 0.03215(66)  & 0.3309(4)   & 1 & 0.00174(31)  \\
  O$_k$ & $4e$ & 0.9597(4) & 0.5339(7) & 0.1689(4)   & 1 & 0.00174(31)  \\
 V & $4e$ & 0.2997(49) & 0.5326(85) & 0.1867(47) & 1 & 0.00127  \\

 \hline \hline
   
\end{tabular}
\end{adjustbox}
    \caption{Averievite, \ave. Atomic positions and displacement parameters from a Rietveld refinement in the $P12_1/c1$ (no. 14) space group using data collected on HRPD (backscattering bank) at $T=1.5$~K. The isotropic displacement parameters were refined for all atoms except for vanadium. }
    \label{avv_1K_params}
\end{table}

\begin{table}[b]
    \centering
    \begin{tabular}{l l l} \hline \hline
     \multicolumn{2}{l}{Cu - Cu}  &  Distance (\AA) \\ \hline \\[-1em] 
     \multicolumn{2}{l}{Cu$_{k1}$ - Cu$_{{k2}(1)}$} & 3.151(4) \\ \\[-1em] 
     \multicolumn{2}{l}{Cu$_{k1}$ - Cu$_{{k2}(2)}$} & 3.225(4) \\ \\[-1em] 
     \multicolumn{2}{l}{Cu$_{{k2}(1)}$ - Cu$_{{k2}(2)}$} & 3.2251(10) \\ \\[-1em]
     \multicolumn{2}{l}{Cu$_{k1}$ - Cu$_h$} & 2.919(4) \\ \\[-1em] 
     \multicolumn{2}{l}{Cu$_{{k2}(1)}$ - Cu$_h$} & 2.913(4) \\ \\[-1em]
     \multicolumn{2}{l}{Cu$_{{k2}(2)}$ - Cu$_h$} & 2.935(4) \\ \\[-1em] 
  
    \hline \hline \\ \hline \hline
     \multicolumn{2}{l}{Cu$_{k2}$ - Cu$_{k2}$} & \multirow{2}{*}{Angle}  \\ \\[-1em] 
     \multicolumn{2}{l}{Cu$_{k1}$ - Cu$_{k2}$} &  \\ \hline \\[-1em] 
     \multicolumn{2}{l}{$\angle$ Cu$_{{k2}(1)}$ - O$_k$ - Cu$_{{k2}(2)}$} & 118.9(3)\degree  \\ \\[-1em] 
     \multicolumn{2}{l}{$\angle$ Cu$_{k1}$ - O$_k$ - Cu$_{{k2}(1)}$} & 114.7(3)\degree \\ \\[-1em] 
     \multicolumn{2}{l}{$\angle$ Cu$_{k1}$ - O$_k$ - Cu$_{{k2}(2)}$} & 117.0(3)\degree  \\ \\[-1em] 
      %
       \hline \hline \\ \hline \hline
       Cu$_{k2}$ - Cu$_h$  & \multirow{2}{*}{Angle} & \multirow{2}{*}{Average}\\ \\[-1em] 
       Cu$_{k1}$ - Cu$_h$ & & \\  \hline \\[-1em] 
       $\angle$ Cu$_{{k2}(1)}$ - O$_k$ - Cu$_h$ & 100.8(3)\degree & \multirow{2}{*}{96.04\degree} \\ \\[-1em] 
       $\angle$ Cu$_{{k2}(1)}$ - O$_1$ - Cu$_h$ & 91.27(18)\degree & \\ \\[-1em] 
       $\angle$ Cu$_{{k2}(2)}$ - O$_k$ - Cu$_h$ & 100.4(3)\degree & \multirow{2}{*}{96.68\degree} \\ \\[-1em] 
       $\angle$ Cu$_{{k2}(2)}$ - O$_2$ - Cu$_h$ & 92.96(19)\degree & \\ \\[-1em] 
       $\angle$ Cu$_{k1}$ - O$_k$ - Cu$_h$ & 99.9(3)\degree & \multirow{2}{*}{95.55\degree} \\ \\[-1em] 
       $\angle$ Cu$_{k1}$ - O$_{Cs}$ - Cu$_h$ & 91.19(17)\degree &  \\ \\[-1em] 

%     
     \hline \hline
%  
\end{tabular}
    \caption{Averievite, \ave, Cu-Cu distances in the tetrahedra of the pyrochlore slabs and Cu-O-Cu superexchange angles, from crystallographic structure refinements in $P12_1/c1$ using data collected at $T=1.5$~K on HRPD. From each Cu$_{k1}$ and Cu$_{k2}$ there are two possible superexchange pathways to Cu$_h$, so the average angle was calculated. %The pathways within the kagome triangles are likely to favour antiferromagnetic exchange interactions, whereas those to the interlayer Cu$_h$ are near the 95$\degree$ crossover from ferro- to antiferromagnetic indicating the exchanges are likely close to zero.
    }
    \label{TableAveAngles}
\end{table}

\clearpage
\subsection{\znave}

Neutron diffraction data collected on HRPD at $T=1.5$~K with a Rietveld refinement of the \znave\ structure in the $P3$ space group is shown in Fig.~\ref{FigZnCu4HRPDRietveld}. The resulting lattice parameters and atomic positions are displayed in Table~\ref{Znavv_1K_params}.

% Figure Rietveld refinement ZnCu4 HRPD 1.5 K
\begin{figure*}[!hb]
\centering
\includegraphics[width=0.8\textwidth]{Cu4_hrpd_1K_sg143_lowQ.pdf}
\includegraphics[width=0.8\textwidth]{Cu4_hrpd_1K_sg143_highQ.pdf}
\caption{\znave\ Rietveld refinement in the $P3$ space group (no. 143) on data collected at 1.5~K on the back-scattering bank of HRPD with 70 parameters ($r_{wp} = 2.78$ and $\chi^2 = 1.96$). There is a $\sim2$\% CuO impurity shown by the blue ticks. The regions not included in the refinement are scattering from the sample can. ({\bf{Top}}) Low-Q range and ({\bf{Bottom}}) high-Q range.
}
\label{FigZnCu4HRPDRietveld}
\end{figure*}
% Source: Madeleine's thesis

\begin{sidewaystable}
    \centering
    \begin{tabular}{ p{2cm} p{1.8cm} p{2cm} p{2.cm} p{2.cm} p{3.cm} p{2.cm}} \hline \hline
    \multicolumn{7}{c}{Lattice parameters} \\ \hline
     \multicolumn{2}{c}{$a$ (\AA)}  & $b$ (\AA) & $c$ (\AA) & $\alpha$ ($\degree$) & $\beta$ ($\degree$) & $\gamma$ ($\degree$)  \\ \hline
     \multicolumn{2}{c}{6.27989(4)} & 6.27989(4) & 8.41826(4) & 90 & 90 & 120  \\ \hline
    \multicolumn{7}{c}{Atomic parameters} \\ \hline
     Atom  & Wyckoff & $x$ & $y$ & $z$ & Occupation & $U_\mathrm{iso}$ (\AA$^2$) \\
     & site & & & & & \\ \hline
 Cu$_{k}$ & $3d$ &  0.5053(55) & 0 & 0  & 1 & 0.06967 \\
 Zn$_{h1}$/Cu$_{h1}$ & $1b$ & 1/3 & 2/3  & 0.7321(3)  & 1 & 0.0125  \\
 Zn$_{h2}$/Cu$_{h2}$ & $1c$ & 2/3 & 1/3  & 0.2678(1)  & 1 & 0.0125  \\
 Cl & $1a$ & 0  & 0 & 0 & 1 & 0.0195 \\
 Cs$_1$ &  $1a$ &  0 &  0 &  0.4586(5) & 0.5 & 0.0354 \\
 Cs$_2$ &  $1a$ &  0 &  0 &  0.5414(6) & 0.5 & 0.0354 \\
 O$_{Cs1}$ & $1b$ & 1/3 & 2/3 & 0.500(37) & 1 & 0.0279  \\
 O$_{Cs2}$ & $1c$ & 2/3 & 1/3 & 0.500(37) & 1 & 0.0279  \\
 O$_1$  & $3d$ & 0.4880(29)  & 0.5255(44)  & 0.2386(1) & 1 & 0.0501 \\
 O$_2$  & $3d$ & 0.4880(29)  & 0.4493(43)  & 0.7613(4) & 1 & 0.0501 \\
 O$_{k1}$  & $1b$ & 1/3  & 2/3  & 0.9586(2) & 1 & 0.0279 \\
  O$_{k2}$  & $1c$ & 2/3  & 1/3  & 0.04141(1) & 1 & 0.0279 \\
 V$_1$  & $1b$ & 1/3  & 2/3  & 0.3061 & 1 & 0.0465 \\
 V$_2$  & $1c$ & 2/3  & 1/3  & 0.6939 & 1 & 0.0465 \\
    \hline 
    \multicolumn{7}{c}{Anisotropic displacement parameters (\AA$^2$)} \\ \hline 
  Atom &  U$_{11}$ & U$_{22}$ & U$_{33}$ & U$_{12}$ & U$_{13}$ & U$_{23}$ \\ \hline 
Cu$_{k}$ & 0.1664(51) & 0.0182(12) & 0.0088(8) & 0.0268(36) &  \minus 0.0332(28) & \minus 0.0055(10) \\
Zn$_{h1}$/Cu$_{h1}$ & 0.1844(60) & 0.1844(60) & 0.0009(9) & 0.0092(3) & 0 & 0 \\
Zn$_{h2}$/Cu$_{h2}$ & 0.1844(60) & 0.1844(60) & 0.0009(9) & 0.0092(3) & 0 & 0 \\
Cl & 0.0431(24) & 0.0431(24) & 0.0203(11) & 0.0414(28) & 0 & 0 \\
Cs$_1$ &  0.0570(30) & 0.0570(30) & 0.0286(32) & 0.0285(15) & 0 & 0 \\
Cs$_2$ &  0.0570(30) & 0.0570(30) & 0.0286(32) & 0.0285(15) & 0 & 0 \\
O$_{Cs1}$ & 0.0254(12) & 0.0537(91) & 0.0059(9) & 0.0215(28) & 0 & 0  \\
O$_{Cs2}$ & 0.0254(12) & 0.0537(91) & 0.0059(9) & 0.0215(28) & 0 & 0  \\ 
O$_1$ & 0.0798(24) & 0.1153(49) & 0.0090(7) & 0.0892(25) & 0.0086(2) & 0.0110(22)  \\
O$_2$ & 0.0798(24) & 0.1153(49) & 0.0090(7) & 0.0892(25) & 0.0086(2) & 0.0110(22)  \\
O$_{k1}$ & 0.0254(12) & 0.0537(91) & 0.0059(9) & 0.0215(28) & 0 & 0  \\
O$_{k1}$ & 0.0254(12) & 0.0537(91) & 0.0059(9) & 0.0215(28) & 0 & 0  \\
V$_1$ & 0.0223 & 0.1079  & 0.0061 &  0.0302   & 0 & 0 \\
V$_2$    & 0.0223 & 0.1079  & 0.0061 &  0.0302   & 0 & 0 
 \\ \hline \hline
   
\end{tabular}
    \caption{\znave. Atomic positions and displacement parameters from the Rietveld refinement in the space group $P3$ (no. 143) using neutron data collected on HRPD (backscattering bank) at $T=1.5$~K. Anisotropic displacement parameters (ADPs) were stably refined for all atoms apart from vanadium. The positions and ADPs of the vanadium atoms were fixed to the values from the synchrotron data at $T=100$~K .}
    \label{Znavv_1K_params}
\end{sidewaystable}

\clearpage
\section{\zntwoave\ crystal structure}

Neutron diffraction data collected on HRPD at $T=1.5$~K with a Rietveld refinement of the \zntwoave\ structure in the $P3$ space group is shown in Fig.~\ref{Zn2avv_1K_hrpd}. 
The resulting lattice parameters and atomic positions are displayed in Table~\ref{Zn2avv_1K_params}.

\begin{figure}[h]
    \centering
    \includegraphics[width=0.8\textwidth]{Cu3_hrpd_1K_sg143_004_peak_stephens_CircleConvolution_hklOffset_140323_lowQ.pdf}
    \includegraphics[width=0.8\textwidth]{Cu3_hrpd_1K_sg143_004_peak_stephens_CircleConvolution_hklOffset_140323_highQ.pdf}
    \caption{\zntwoave\ Rietveld refinement in the $P3$ space group (no. 143) on data collected at $T=1.5$~K on the back-scattering bank of HRPD with 82 parameters ($r_{wp} = 3.12$ and $\chi^2 = 3.15$). There is a $\sim 2$\% CuO impurity shown by the blue ticks and unidentified impurity peaks are marked by asterisks. The regions not included in the refinement are scattering from the aluminium sample can. \textbf{Top} Low-$Q$ range and \textbf{Bottom} high-$Q$ range.}
    \label{Zn2avv_1K_hrpd}
         
\end{figure}


\clearpage
\begin{sidewaystable}
    \centering
    \begin{tabular}{ p{2cm} p{1.8cm} p{2cm} p{2.cm} p{2.cm} p{3.cm} p{2.cm}} \hline \hline
    \multicolumn{7}{c}{Lattice parameters} \\ \hline
     \multicolumn{2}{c}{$a$ (\AA)} & $b$ (\AA) & $c$ (\AA) & $\alpha$ ($\degree$) & $\beta$ ($\degree$) & $\gamma$ ($\degree$)    \\ \hline
     \multicolumn{2}{c}{6.24863(4)} & 6.24863(4) & 8.48754(8) & 90 & 90 & 120 \\ \hline
    \multicolumn{7}{c}{Atomic parameters} \\ \hline
     Atom  & Wyckoff site & $x$ & $y$ & $z$ & Occupation & $U_\mathrm{iso}$ (\AA$^2$) \\ \hline
 Cu$_k$/Zn$_k$ & $3d$ &  0.523(2) & 0 & 0  & 0.95(4)/0.05(4)& 0.0784 \\
 Zn$_{h1}$/Cu$_{h1}$ & $1b$ & 1/3 & 2/3  & 0.7307(3)  & 0.93(3)/0.07(3) & 0.0287  \\
 Zn$_{h2}$/Cu$_{h2}$ & $1c$ & 2/3 & 1/3  & 0.2693(3)  & 0.93(3)/0.07(3) & 0.0287  \\
 Cl & $1a$ & 0  & 0 & 0 & 1 & 0.0246  \\
 Cs1 &  $1a$ &  0 &  0 &  0.480(6) & 1 & 0.0390 \\
 Cs2 &  $1a$ &  0 &  0 &  0.520(6) & 1 & 0.0390 \\
 O$_{Cs1}$ & $1b$ & 1/3 & 2/3 & 0.502(2) & 1 & 0.0246  \\
 O$_{Cs2}$ & $1c$ & 2/3 & 1/3 & 0.502(2) & 1 & 0.0246  \\
 O$_{h1}$  & $3d$ & 0.491(3) & 0.520(4) & 0.2400(2) & 1 & 0.0988 \\
 O$_{h2}$  & $3d$ & 0.491(3) & 0.447(4) & 0.7600(2) & 1 & 0.0988 \\
 O$_{k1}$  & $1b$ & 1/3  & 2/3  & 0.9609(3) & 1 & 0.0427 \\
 O$_{k2}$  & $1c$ & 2/3  & 1/3  & 0.0391(3) & 1 & 0.0427 \\
 V$_1$  & $1b$ & 1/3  & 2/3  & 0.30743 & 1 & 0.0379 \\
 V$_2$  & $1c$ & 2/3  & 1/3  & 0.69257 & 1 & 0.0379 \\
    \hline 
    \multicolumn{7}{c}{Anisotropic displacement parameters (\AA$^2$)} \\ \hline 
  Atom &  U$_{11}$ & U$_{22}$ & U$_{33}$ & U$_{12}$ & U$_{13}$ & U$_{23}$ \\ \hline 
Cu$_k$/Zn$_k$ & 0.189(9) & 0.0280(2) & 0.018(2) & 0.030(5) &  -0.050(42) & -0.006(2) \\
Zn$_{h1}$/Cu$_{h1}$ & 0.020(2) & 0.06(1) & 0.004(2) & 0.031(5) & 0 & 0 \\
Zn$_{h2}$/Cu$_{h2}$ & 0.020(2) & 0.06(1) & 0.004(2) & 0.031(5) & 0 & 0 \\
Cl & 0.027(3) & 0.027(3) & 0.0193(16) & 0.020(3) & 0 & 0 \\
Cs1 &  0.007(3) & 0.01(4) & 0.06(2) & -0.007(12)  & 0 & 0 \\
Cs2 &  0.007(3) & 0.01(4) & 0.06(2) & -0.007(12)  & 0 & 0 \\
O$_{Cs1}$ & 0.0053(17) & 0.07(3) & 0.0086(13) & 0.007(8) & 0 & 0 \\
O$_{Cs2}$ & 0.0053(17) & 0.07(3) & 0.0086(13) & 0.007(8) & 0 & 0 \\
O$_{h1}$ & 0.121(3) & 0.162(5) & 0.0124(11) & 0.135(3) & 0.025(3) & 0.029(3) \\
O$_{h2}$ &  0.121(3) & 0.162(5) & 0.0124(11) & 0.135(3) & 0.025(3) & 0.029(3) \\
O$_{k1}$ & 0.034(4) & 0.08(2) & 0.0063(16) & 0.035(8) & 0 & 0 \\
O$_{k2}$ & 0.034(4) & 0.08(2) & 0.0063(16) & 0.035(8) & 0 & 0 \\
V$_{1}$ & 0.0065 & 0.0987 & 0.0085 & 0.0102 & 0 & 0  \\
V$_{2}$ & 0.0065 & 0.0987 & 0.0085 & 0.0102 & 0 & 0  
 \\ \hline \hline
   
\end{tabular}
    \caption{\zntwoave. Atomic positions and displacement parameters from the Rietveld refinement in the $P3$ space group (no. 143) using neutron data collected on HRPD (backscattering bank) at $T=1.5$~K. Anisotropic displacement parameters (ADPs) were stably refined for all atoms apart from vanadium. The position and ADPs of vanadium were fixed at the values refined from the synchrotron data at $T=100$~K. The refined values of the Cu/Zn site occupancies indicate a sample with nominal stoichiometry and minimal antisite disorder.}
    \label{Zn2avv_1K_params}

\end{sidewaystable}

\clearpage
\section{Scanning electron microscopy (SEM)}

SEM images of \znave\ and \zntwoave\ indicating platelets with no amorphous impurity phase in either sample.

% Figure SEM images of ZnCu4 and Zn2Cu3
\begin{figure*}[!h]
\includegraphics[width=0.49\textwidth]{MADELEINE_S1_003_png.png}
\includegraphics[width=0.49\textwidth]{MADELEINE_002_png.png}
\caption{SEM images of ({\bf left}) \znave\ showing platelets of various sizes up to a few \textmu m and no amorphous impurity phase and ({\bf right}) \zntwoave\ showing 2-dimensional crystallites up to $\sim1$~\textmu m in diameter and no amorphous impurity phase.
}
\label{FigSEM}
\end{figure*}
% Source: Madeleine's thesis

\section{Magnetometry}

DC susceptibility measurements for \znave\ showing an anomaly at $T_1=3.5$~K.

% Figure Susceptibility and first derivative of ZnCu4 

\begin{figure*}[!h]
\includegraphics[width=0.49\textwidth]{ZnCu4_2021_ZFC_FC.pdf}
\includegraphics[width=0.49\textwidth]{ZnCu4_2021_chi_derivative.pdf}
\caption{\znave . ({\bf Left}) Magnetic susceptibility collected in a field of 0.1~T. ({\bf Right}) First derivative of magnetic susceptibility showing an anomaly at $T_1=3.5$~K.
}
\label{FigZnCu4Susceptibility}
\end{figure*}

\clearpage
\section{Magnetic structure of averievite}

The basis vectors for the irreducible representations $\Gamma_1$ and $\Gamma_3$ in space group \textit{P}12$_\mathbf{1}$/\textit{c}1 with \textit{k}=(0.5,~0,~0) are given in Tables \ref{avv_cuk1_BVs_gamma1}-\ref{avv_cuh_BVs_gamma3} and the refined mixing coefficients for each basis vector are given in Table~\ref{avv_refined_BV_coeffs}. The magnetic structure refinement in $\Gamma_3$ of the $P12_1/c1$ space group using neutron diffraction data collected on WISH is shown in Fig.~\ref{FigCu5MagStrRefinement}. 

\begin{table*}[ht!]
    \centering
    \begin{tabular}{c c c c c c} \hline \hline
      Atom number & Coordinates & {Basis vector} & {m$_a$} & m$_b$ &{m$_c$} \\ \hline
       Atom 1 & (0,~0,~0.5) & $\psi_1$ & 1 & 0 & 0 \\ 
      & & $\psi_2$ & 0 & 1 &  0  \\
     & & $\psi_3$ & 0 & 0 & 1  \\
     Atom 2 & (0,~0.5,~0) & $\psi_1$ &  \minus 1 & 0 & 0 \\
    & & $\psi_2$ & 0 &  1 &  0  \\
     & & $\psi_3$ & 0 & 0 &  \minus 1  \\ 
       \hline \hline
   
\end{tabular}
    \caption{The basis vectors of the $\Gamma_1$ irreducible representation of the space group $P12_1/c1$ with ${k}=(0.5,~0,~0)$ for the Cu$_{k1}$ $2c$ site, where m$_a$, m$_b$ and m$_c$ are the moment components along the crystallographic unit cell directions $a$, $b$ and $c$.}
    \label{avv_cuk1_BVs_gamma1}
\end{table*}

\begin{table*}[ht!]
    \centering
    \begin{tabular}{c c c c c c} \hline \hline
      Atom number & Coordinates & {Basis vector} & {m$_a$} & m$_b$ &{m$_c$} \\ \hline
       Atom 1 & (0,~0.80048,~0.23059) & $\psi_1$ & 1 & 0 & 0 \\ 
      & & $\psi_2$ & 0 & 1 &  0  \\
     & & $\psi_3$ & 0 & 0 & 1  \\
     Atom 2 & (0,~0.30048,~0.26941) & $\psi_1$ & \minus 1 & 0 & 0 \\
    & & $\psi_2$ & 0 &  1 &  0  \\
     & & $\psi_3$ & 0 & 0 & \minus 1  \\ 
     Atom 3 & (0,~0.19952,~0.76941) & $\psi_1$ &  1 & 0 & 0 \\
    & & $\psi_2$ & 0 &  1 &  0  \\
     & & $\psi_3$ & 0 & 0 &  1  \\ 
     Atom 4 & (0,~0.69952,~0.73059) & $\psi_1$ & \minus 1 & 0 & 0 \\
    & & $\psi_2$ & 0 &  1 &  0  \\
     & & $\psi_3$ & 0 & 0 & \minus 1 \\ 
       \hline \hline
   
\end{tabular}
    \caption{The basis vectors of the $\Gamma_1$ irreducible representation of the space group $P12_1/c1$ with ${k}=(0.5,~0,~0)$ for the Cu$_{k2}$ $4e$ site, where m$_a$, m$_b$ and m$_c$ are the moment components along the crystallographic unit cell directions $a$, $b$ and $c$.}
    \label{avv_cuk2_BVs_gamma1}
\end{table*}

\begin{table*}[ht!]
    \centering
    \begin{tabular}{c c c c c c} \hline \hline
      Atom number & Coordinates & {Basis vector} & {m$_a$} & m$_b$ &{m$_c$} \\ \hline
       Atom 1 & (0.72862,~0.53568,~0.16828) & $\psi_1$ & 1 & 0 & 0 \\ 
      & & $\psi_2$ & 0 & 1 &  0  \\
     & & $\psi_3$ & 0 & 0 & 1  \\
     Atom 2 & (0.27138,~0.03568,~0.33172) & $\psi_1$ & 1 & 0 & 0 \\
    & & $\psi_2$ & 0 & \minus 1 &  0  \\
     & & $\psi_3$ & 0 & 0 & 1  \\ 
     Atom 3 & (0.27138,~0.46432,~0.83172) & $\psi_1$ & \minus 1 & 0 & 0 \\
    & & $\psi_2$ & 0 & \minus 1 &  0  \\
     & & $\psi_3$ & 0 & 0 & \minus 1  \\ 
     Atom 4 & (0.72862,~0.96432,~0.66828) & $\psi_1$ & \minus 1 & 0 & 0 \\
    & & $\psi_2$ & 0 &  1 &  0  \\
     & & $\psi_3$ & 0 & 0 & \minus 1 \\ 
       \hline \hline
   
\end{tabular}
    \caption{The basis vectors of the $\Gamma_1$ irreducible representation of the space group $P12_1/c1$ with ${k}=(0.5,~0,~0)$ for the Cu$_{h}$ $4e$ site, where m$_a$, m$_b$ and m$_c$ are the moment components along the crystallographic unit cell directions $a$, $b$ and $c$.}
    \label{avv_cuh_BVs_gamma1}
\end{table*}

\begin{table*}[ht!]
    \centering
    \begin{tabular}{c c c c c c} \hline \hline
      Atom number & Coordinates & {Basis vector} & {m$_a$} & m$_b$ &{m$_c$} \\ \hline
       Atom 1 & (0,~0,~0.5) & $\psi_1$ & 1 & 0 & 0 \\ 
      & & $\psi_2$ & 0 & 1 &  0  \\
     & & $\psi_3$ & 0 & 0 & 1  \\
     Atom 2 & (0,~0.5,~0) & $\psi_1$ &  1 & 0 & 0 \\
    & & $\psi_2$ & 0 &  \minus 1 &  0  \\
     & & $\psi_3$ & 0 & 0 & 1  \\ 
       \hline \hline
   
\end{tabular}
    \caption{The basis vectors of the $\Gamma_3$ irreducible representation of the space group $P12_1/c1$ with ${k}=(0.5,~0,~0)$ for the Cu$_{k1}$ $2c$ site, where m$_a$, m$_b$ and m$_c$ are the moment components along the crystallographic unit cell directions $a$, $b$ and $c$.}
    \label{avv_cuk1_BVs_gamma3}
\end{table*}

\begin{table*}[ht!]
    \centering
    \begin{tabular}{c c c c c c} \hline \hline
      Atom number & Coordinates & {Basis vector} & {m$_a$} & m$_b$ &{m$_c$} \\ \hline
       Atom 1 & (0,~0.80048,~0.23059) & $\psi_1$ & 1 & 0 & 0 \\ 
      & & $\psi_2$ & 0 & 1 &  0  \\
     & & $\psi_3$ & 0 & 0 & 1  \\
     Atom 2 & (0,~0.30048,~0.26941) & $\psi_1$ &  1 & 0 & 0 \\
    & & $\psi_2$ & 0 &  \minus 1 &  0  \\
     & & $\psi_3$ & 0 & 0 &  1  \\ 
     Atom 3 & (0,~0.19952,~0.76941) & $\psi_1$ &  1 & 0 & 0 \\
    & & $\psi_2$ & 0 &  1 &  0  \\
     & & $\psi_3$ & 0 & 0 &  1  \\ 
     Atom 4 & (0,~0.69952,~0.73059) & $\psi_1$ &  1 & 0 & 0 \\
    & & $\psi_2$ & 0 & \minus 1 &  0  \\
     & & $\psi_3$ & 0 & 0 &  1 \\ 
       \hline \hline
   
\end{tabular}
    \caption{The basis vectors of the $\Gamma_3$ irreducible representation of the space group $P12_1/c1$ with ${k}=(0.5,~0,~0)$ for the Cu$_{k2}$ $4e$ site, where m$_a$, m$_b$ and m$_c$ are the moment components along the crystallographic unit cell directions $a$, $b$ and $c$.}
    \label{avv_cuk2_BVs_gamma3}
\end{table*}

\begin{table*}[ht!]
    \centering
    \begin{tabular}{c c c c c c} \hline \hline
      Atom number & Coordinates & {Basis vector} & {m$_a$} & m$_b$ &{m$_c$} \\ \hline
       Atom 1 & (0.72862,~0.53568,~0.16828) & $\psi_1$ & 1 & 0 & 0 \\ 
      & & $\psi_2$ & 0 & 1 &  0  \\
     & & $\psi_3$ & 0 & 0 & 1  \\
     Atom 2 & (0.27138,~0.03568,~0.33172) & $\psi_1$ & \minus 1 & 0 & 0 \\
    & & $\psi_2$ & 0 &  1 &  0  \\
     & & $\psi_3$ & 0 & 0 & \minus 1  \\ 
     Atom 3 & (0.27138,~0.46432,~0.83172) & $\psi_1$ & \minus 1 & 0 & 0 \\
    & & $\psi_2$ & 0 & \minus 1 &  0  \\
     & & $\psi_3$ & 0 & 0 & \minus 1  \\ 
     Atom 4 & (0.72862,~0.96432,~0.66828) & $\psi_1$ &  1 & 0 & 0 \\
    & & $\psi_2$ & 0 &  \minus 1 &  0  \\
     & & $\psi_3$ & 0 & 0 &  1 \\ 
       \hline \hline
   
\end{tabular}
    \caption{The basis vectors of the $\Gamma_3$ irreducible representation of the space group $P12_1/c1$ with ${k}=(0.5,~0,~0)$ for the Cu$_{h}$ $4e$ site, where m$_a$, m$_b$ and m$_c$ are the moment components along the crystallographic unit cell directions $a$, $b$ and $c$.}
    \label{avv_cuh_BVs_gamma3}
\end{table*}

\begin{table*}[htb!]
    \centering
    \begin{tabular}{c c c c} \hline \hline
     Atom & Coordinates & Basis vector, $\psi_i$ & $C_{i}$  \\ \hline
     Cu$_{k1}$ & (0,~0,~0.5) & $\psi_1$ & 0\\
               & & $\psi_2$ & 0 \\
               & & $\psi_3$ & \minus0.30(1) \\ 
    
    Cu$_{k2}$ & (0,~0.80048,~0.23059) & $\psi_1$ & 0 \\
              & & $\psi_2$ & \minus0.33(1) \\
              & & $\psi_3$ & \minus0.33(1) \\
                     
    Cu$_h$ & (0.72862,~0.53568,~0.16828) & $\psi_1$ & 0\\
           & & $\psi_2$ & \minus0.215(8) \\
           & & $\psi_3$ & \minus0.56(1)\\                 

    \hline \hline
\end{tabular}
    \caption{The refined mixing coefficients, $C_{i}$, for each basis vector of Atom 1 of the three Cu sites, Cu$_{k1}$, Cu$_{k2}$ and Cu$_h$, (see Tables~\ref{avv_cuk1_BVs_gamma1}-\ref{avv_cuh_BVs_gamma3}) corresponding to the refinement shown in Fig.~\ref{FigCu5MagStrRefinement}.  }
    \label{avv_refined_BV_coeffs}
\end{table*}

% Figure magnetic structure refinement in Gamma 3

\begin{figure*}[!h]
\centering
\includegraphics[width=0.99\textwidth]{avv_final_Gamma3_refinement_bank1.pdf}
\includegraphics[width=0.99\textwidth]{avv_final_Gamma3_refinement_bank2.pdf}
\caption{Averievite WISH data obtained from the temperature subtraction $T = 1.5-40$~K, with a magnetic structure refinement (red) in $\Gamma_3$. The grey regions arise from the subtraction of nuclear peaks and were excluded from the refinements. The green tick marks indicate the magnetic Bragg peak positions and the difference plot is shown in black. ({\bf Top}) Bank 1. ({\bf Bottom}) Bank 2.
}
\label{FigCu5MagStrRefinement}
\end{figure*}


\end{document}